\documentclass[../../main.tex]{subfiles}
\begin{document}

{\bf [解]}
簡単のため$s=\sin x,\ c=\cos x$と書く．
まずは$y=\tan x$と$y=\cos x$の交点の$x$座標を求める．$y$を消去した方程式をとくと
\begin{align}
\tan x=\cos x \quad\therefore\ s=c^2=1-s^2 \quad\therefore\ s=\dfrac{1}{2}(-1\pm\sqrt{5})
\end{align}
である．このうち$|s|\le 1$を満たすのは
\begin{align}
 s=\dfrac{-1+\sqrt{5}}{2}\equiv s_0 \label{eq:1}
\end{align}
である．そこで$\sin x=s_0$を満たす$x$を$\alpha$とおく．$|\alpha|<\dfrac{\pi}{2}$より$\cos\alpha>0$だから
\begin{align}
 \cos\alpha 
  &= \sqrt{1-s^2_0} \\
  &= \sqrt{1-\left(\dfrac{-1+\sqrt{5}}{2}\right)^2} \\ 
  &= \sqrt{\dfrac{-1+\sqrt{5}}{2}} \label{eq:2}
\end{align}
である．
$(\tan x)'=\dfrac{1}{c^2}$だから\cref{eq:1}と合わせて接線$\ell$の傾きは
\begin{align}
\dfrac{1}{1-s_0^2}=\dfrac{2}{\sqrt{5}-1}=\dfrac{\sqrt{5}+1}{2} \label{eq:3}
\end{align}
である．
$A$から$x$軸に下ろした垂足$B$，$\ell$と$x$軸の交点$C$とするともとめる面積$T$として
\begin{align}
  T=\int_0^\alpha\tan x\,dx-\triangle ABC \quad\cdots\text{①} \label{eq:4}
\end{align}

\begin{figure}[htb]
\begin{center}
\begin{tikzpicture}[scale=3.0, >=stealth, every node/.style={font=\small}]

  % 1. 数値・座標の設定
  % \alpha = arcsin((-1+sqrt(5))/2) \approx 0.6662 rad (38.17°)
  % y_A = tan(\alpha) = cos(\alpha) = sqrt((-1+sqrt(5))/2) \approx 0.7862
  % 接線 l: y = 1.6180 * (x - 0.6662) + 0.7862
  % 接線と x 軸の交点 C の x 座標: 0.6662 - 0.7862 / 1.6180 \approx 0.1802

  \coordinate (A) at (0.6662, 0.7862);
  \coordinate (B) at (0.6662, 0);
  \coordinate (C) at (0.1802, 0);

  % 2. 面積 T の斜線塗りつぶし
  % 領域 T : x = 0 から x = \alpha(0.6662) までの y = tan x の下側領域から △ABC を差し引いた部分
  \fill[pattern=north east lines, pattern color=black!50]
    (0,0) -- plot[domain=0:0.6662, samples=50] (\x, {tan(\x r)})
    -- (C) -- (0,0) -- cycle;

  % 3. 補助点線 (A から B への垂線)
  \draw[dashed, gray!70, thick] (A) -- (B);

  % 4. 座標軸
  \draw[->, thick] (-0.6, 0) -- (1.3, 0) node[right] {$x$};
  \draw[->, thick] (0, -0.3) -- (0, 1.6) node[above] {$y$};
  \node[above left] at (0,0) {$O$};

  % 5. 曲線・直線の描画
  \begin{scope}
    \clip (-0.6, -0.3) rectangle (1.2, 1.5);
    % y = tan x
    \draw[domain=-0.6:0.95, smooth, thick, blue!80!black, samples=60] 
      plot (\x, {tan(\x r)});
    % y = cos x
    \draw[domain=-0.6:1.1, smooth, thick, green!50!black, samples=60] 
      plot (\x, {cos(\x r)});
    % 接線 l (点 C から点 A を通る直線)
    \draw[domain=-0.1:1.0, thick, red!80!black] 
      plot (\x, {1.6180*(\x - 0.6662) + 0.7862});
  \end{scope}

  % 6. 曲線・直線の名札ラベル
  \node[blue!80!black, left] at (0.85, 1.25) {$y=\tan x$};
  \node[green!50!black, right] at (0.85, 0.55) {$y=\cos x$};
  \node[red!80!black, right] at (0.85, 1.1) {$\ell$};

  % 7. 特徴的なプロット点とラベル
  \fill (A) circle (1.0pt) node[above left] {$A$};
  \fill (B) circle (1.0pt) node[below] {$B$};
  \fill (C) circle (1.0pt) node[below] {$C$};

  % 8. 領域 T の名札
  \node[fill=white, inner sep=1pt, font=\bfseries] at (0.32, 0.22) {$T$};

\end{tikzpicture}
\end{center}
\caption{曲線 $y=\tan x$，$y=\cos x$ と面積 $T$（斜線部）}
\end{figure}


まず$\triangle ABC$の面積は\cref{eq:2,eq:3}より
\begin{align}
\triangle ABC
  &=\dfrac{1}{2}|AB||AC| \\
  &=\dfrac{1}{2}|AB|\dfrac{|AB|}{\dfrac{\sqrt{5}+1}{2}} \\
  &=\dfrac{1}{2}\dfrac{|\cos\alpha|^2}{\dfrac{\sqrt{5}+1}{2}} \\
  &=\dfrac{1}{2}\left(\sqrt{\dfrac{\sqrt{5}-1}{2}}\right)^2\cdot\dfrac{1}{\dfrac{\sqrt{5}+1}{2}} \\
  &=\dfrac{1}{2}\cdot\dfrac{\sqrt{5}-1}{\sqrt{5}+1} \\
  &=\dfrac{1}{2}\cdot\dfrac{(\sqrt{5}-1)^2}{4}\quad\cdots\text{②} \label{eq:5}
\end{align}
である．

\cref{eq:4}の第一項は
\begin{align}
\int_0^\alpha\tan x\,dx
  &=\left[-\log(\cos x)\right]_0^\alpha \\
  &=-\log(\cos\alpha) \\
  &=-\dfrac{1}{2}\log\dfrac{\sqrt{5}-1}{2}\quad\cdots\text{③} \label{eq:6}
\end{align}
である．

\cref{eq:5,eq:6}を\cref{eq:4}に代入して求める面積は
\begin{align}
 T=-\dfrac{1}{2}\log\dfrac{\sqrt{5}-1}{2}-\dfrac{1}{8}(\sqrt{5}-1)^2
\end{align}
である．


\end{document}
